\begingroup
\renewcommand{\arraystretch}{1.45}
\clearpage\fancyhead[L]{\small ARITHMETIC–GEOMETRIC MEAN}
\pdfbookmark[0]{ARITHMETIC–GEOMETRIC MEAN}{work-native370}
\SourceLeaf{370}
\selectlanguage{english}
\begin{center}
\textbf{POSTHUMOUS WRITINGS.}\par
\textbf{[ARITHMETIC–GEOMETRIC MEAN.]}\par
\textbf{PART I.}\par
\textbf{ON THE ORIGIN AND}\\\textbf{GENERAL PROPERTIES OF}\par
\textbf{ARITHMETIC–GEOMETRIC MEANS.}\par
\rule{0.09\linewidth}{0.4pt}\par
\end{center}
\AGMArticle{1}
\Needspace{22\baselineskip}
Let $\left\{\begin{array}{l}a,\allowbreak \ a',\allowbreak \ a'',\allowbreak \ a'''\ .\,\allowbreak .\,\allowbreak .\,\allowbreak .\\ b,\allowbreak \ b',\allowbreak \ b'',\allowbreak \ b'''\ .\,\allowbreak .\,\allowbreak .\,\allowbreak .\end{array}\right\}$ be two sequences of quantities formed by the rule that
their corresponding terms are \textit{means} of the preceding terms,
the terms of the upper sequence being arithmetic means and those of the lower sequence geometric means; thus
\[
\begin{aligned}&a'=\tfrac12(a+b),\quad b'=\surd ab,\\[6pt]&a''=\tfrac12(a'+b'),\quad b''=\surd a'b',\\[6pt]&a'''=\tfrac12(a''+b''),\quad b'''=\surd a''b''\ \mathrm{etc.}\end{aligned}
\]
We shall assume that $a$, $b$ are positive real numbers, and that the positive values
of the square roots are taken throughout; the sequences can then be continued
as far as desired, all their terms will be fully determined, and their values will be positive and real. The following observations immediately present themselves:
\par
I. If $a=b$, all the terms of both sequences will be $=a=b$.
\par
II. If, however, $a$, $b$ are unequal, then $(a'-b')(a'+b')=\tfrac14(a-b)^2$, whence
it follows that $b'<a'$; likewise, $b''<a''$, $b'''<a'''$, etc., that is, every term of the lower sequence will be smaller than the corresponding term of the upper sequence. In this case
we shall therefore also assume $b<a$.
\par
\Needspace{13\baselineskip}
III. Under the same assumption, $a'<a$, $b'>b$; $a''<a'$, $b''>b'$, etc.; thus the upper
sequence steadily decreases and the lower steadily increases. It is therefore evident
that each has a limit; these limits are conveniently denoted by $a^{\infty}$, $b^{\infty}$.
\par
\Needspace{19\baselineskip}
IV. Finally, from $\frac{a'-b'}{a-b}=\frac{(a-b)}{4(a'+b')}=\frac{a-b}{2(a+b)+4b'}$ it follows that \hfill\\{}
$a'-\nobreak{}b'<\nobreak{}\tfrac12(a-\nobreak{}b)$;
likewise, $a''-b''<\tfrac12(a'-b')$, etc. It follows that $a-b$, $a'-b'$,
\SourceLeaf{371}
$a''-b''$, $a'''-b'''$, etc. form a steadily decreasing sequence whose
limit is $=0$. Hence $a^{\infty}=b^{\infty}$, that is, the upper and lower sequences will have the same
limit, the former always remaining greater than it and the latter smaller.
\par
We call this limit the \textit{arithmetic–geometric mean between $a$ and $b$},
and denote it by $\mathrm{M}(a,b)$.
\Needspace{35\baselineskip}
\AGMArticle{2}
The roots of the equation $xx-2ax+bb=0$ will be positive real numbers provided
$a\geqq b$; their arithmetic mean will be $a$ and their geometric mean $b$. Thus, if
one root, namely the larger if they are unequal, is denoted by ${}^{\prime}\!a$, and the other by ${}^{\prime}\!b$, ${}^{\prime}\!a$ can be regarded as the term of the upper sequence preceding $a$, and similarly ${}^{\prime}\!b$ as the term of the lower sequence preceding $b$. Likewise, if the roots are denoted as follows,
\begingroup\fontsize{9}{11}\selectfont
\begingroup\setlength{\arraycolsep}{2pt}
\[
\begin{array}{@{}lllll@{}}
\mathrm{equation}&xx-2\,{}^{\prime}\!ax+{}^{\prime}\!b\,{}^{\prime}\!b=0&\mathrm{larger\ root\ denoted\ by}&{}^{\prime\prime}\!a&\mathrm{smaller\ by}\ {}^{\prime\prime}\!b\\[6pt]
&xx-2\,{}^{\prime\prime}\!ax+{}^{\prime\prime}\!b\,{}^{\prime\prime}\!b=0&&{}^{\prime\prime\prime}\!a&\hfill{}^{\prime\prime\prime}\!b\\[6pt]
&xx-2\,{}^{\prime\prime\prime}\!ax+{}^{\prime\prime\prime}\!b\,{}^{\prime\prime\prime}\!b=0&&{}^{\prime\prime\prime\prime}\!a&\hfill{}^{\prime\prime\prime\prime}\!b
\end{array}
\]\endgroup
\endgroup
${}^{\prime}\!a$, ${}^{\prime\prime}\!a$, ${}^{\prime\prime\prime}\!a$, etc. can be regarded as the continuation of the upper sequence to the left, and ${}^{\prime}\!b$, ${}^{\prime\prime}\!b$, ${}^{\prime\prime\prime}\!b$, etc. as the continuation of the lower sequence. There are thus now
two sequences that can be continued indefinitely in both directions:
\[
\begin{array}{@{}ll@{}}
.\,.\,.\,.\ {}^{\prime\prime\prime\prime}\!a,\;{}^{\prime\prime\prime}\!a,\;{}^{\prime\prime}\!a,\;{}^{\prime}\!a,\; a,\; a',\; a'',\; a''',\; a''''\ .\,.\,.\,.\,.&\mathrm{(I)}\\[6pt]
.\,.\,.\,.\ {}^{\prime\prime\prime\prime}\!b,\;{}^{\prime\prime\prime}\!b,\;{}^{\prime\prime}\!b,\;{}^{\prime}\!b,\; b,\; b',\; b'',\; b''',\; b''''\ .\,.\,.\,.\,.&\mathrm{(II)}
\end{array}
\]
\Needspace{32\baselineskip}
\par Every term of sequence (I) will therefore be greater than the corresponding term of sequence (II); the former steadily decreases from left to right and increases
from right to left; the latter steadily increases from left to right and decreases in the opposite direction.
To the right, both sequences have the same limit; to the left, however, (I)
increases beyond every bound, while (II) has the limit $0$ (unless all the terms of both sequences are equal). For ${}^{\prime}\!a=\nobreak{}a+\nobreak{}\surd(aa-bb)$; ${}^{\prime}\!b=a-\surd(aa-bb)$;
hence ${}^{\prime}\!a\,{}^{\prime}\!a-{}^{\prime}\!b\,{}^{\prime}\!b=4a\surd(aa-bb)>4(aa-bb)$, and likewise ${}^{\prime\prime}\!a\,\allowbreak {}^{\prime\prime}\!a-{}^{\prime\prime}\!b\,\allowbreak {}^{\prime\prime}\!b>4({}^{\prime}\!a\,\allowbreak {}^{\prime}\!a-{}^{\prime}\!b\,\allowbreak {}^{\prime}\!b)$,
etc. It is therefore evident that the sequence $aa-bb$, ${}^{\prime}\!a\,{}^{\prime}\!a-{}^{\prime}\!b\,{}^{\prime}\!b$, ${}^{\prime\prime}\!a\,{}^{\prime\prime}\!a-{}^{\prime\prime}\!b\,{}^{\prime\prime}\!b$, etc., and consequently also
$a$, ${}^{\prime}\!a$, ${}^{\prime\prime}\!a$, etc., can exceed every bound; $\frac{{}^{\prime}\!b}{b}=\frac{b}{{}^{\prime}\!a}<\frac{b}{a}$, and likewise $\frac{{}^{\prime\prime}\!b}{{}^{\prime}\!b}<\frac{{}^{\prime}\!b}{{}^{\prime}\!a}$, etc.,
that is, $\frac{{}^{n+1}\!b}{{}^{n}\!b}$ can be made smaller than any bound by increasing $n$, and consequently the limit
of the sequence $b$, ${}^{\prime}\!b$, ${}^{\prime\prime}\!b$, etc. is $=0$.
\Needspace{9\baselineskip}
\SourceLeaf{372}
\par 
The following theorem is so evident from the definition of the arithmetic–geometric mean
that no further explanation is now needed:

\textsc{Theorem.}\quad \textit{The mean between any corresponding terms of sequences I, II is the same as that between $a$ and $b$.}

\clearpage
\endgroup
